% ===================================================================
%  نقشہ نمبر 6 — گھر دے اندر۔ — ساز سامان۔ — کھاݨا۔ — چاݨن۔
%  Traduction 2026 d'apres texto/io, controlee sur texto/fr, texto/en
%  et texto/hi. Consigne : voir texto/pnb/10-naqsha-01.tex.
%
%  Deux notes, toutes deux marquees « (*) », et deux appels : celui de
%  « babo » au bloc c1-12-1, celui de « lambrequino » au bloc c2-01-1.
%  L'ordre les departage, comme dans les neuf autres colonnes.
%
%  LA FEMME DE CHAMBRE N'A PAS DE GROS PLAN, et c'est normal. L'ido du
%  bloc c1-06-1 porte « la chambristino (6) », le francais « la femme
%  de chambre (16) » : c'est l'ido qui se trompe, son propre numero 6
%  etant le savon du bloc c1-02-1. Le pendjabi suit donc le francais
%  et ecrit (16) ; la planche ne porte pas de 16 releve, et le renvoi
%  reste du texte ordinaire. C'est un des trois ecarts declares dans
%  outils/renvoji.py.
% ===================================================================

%%K t06-apar-1 apar
\VUsaut{6.0mm}
\VUtitre{15.0pt}{60}{نقشہ نمبر 6}
\VUsaut{1.4mm}
\VUtitre{11.4pt}{0}{گھر دے اندر، ساز سامان، کھاݨا،}
\VUsaut{0.4mm}
\VUtitre{11.4pt}{0}{چاݨن۔}
\VUsaut{2.2mm}

%%K t06-apar-2 apar
گھر بݨ کے تیار اے۔ جان دے چاچا اپݨے نویں گھر وچ آ وسے نیں، جیدی بݨتر اسیں جاݨنے
ای آں۔ ہُݨ ویکھیے پئی اوہناں اوہنوں کِویں سجایا اے۔ پہلاں اسیں ویکھاں گے:

%%K t06-tit-1 sub
\VUpk{10.2pt}{40}{سوݨ دا کمرہ۔}
\VUsaut{3.0mm}

%%K t06-c1-01-1 p
1. --- اسیں کجھ بے ویلے آ گئے آں، کیوں جے نوکراݨی \VUgras{کمرے} دی صفائی وچ لگی
اے۔

%%K t06-c1-02-1 p
2. --- \VUgras{ہتھ مونہہ دھوݨ دی چوکی} \textsuperscript{(1)} اُتے ایدھر اودھر اوہ
چیزاں پئیاں نیں جیہناں نال نہایا، کنگھی کیتی تے تیار ہویا جاندا اے۔ اوہ نیں
\VUgras{چلمچی} \textsuperscript{(2)} تے \VUgras{پاݨی دا جگ} \textsuperscript{(3)}،
\VUgras{والاں دا برش} \textsuperscript{(4)}، \VUgras{کنگھی} \textsuperscript{(5)}،
\VUgras{صابن} \textsuperscript{(6)} تے \VUgras{اسفنج} \textsuperscript{(7)}، تے
چھیکر کُلی کرن دے پاݨی دی اک نِکی \VUgras{شیشی} \textsuperscript{(8)}۔

%%K t06-c1-03-1 p
3. --- زمین اُتے، چوکی دے سامݨے، ایہہ رہی گندے پاݨی دی \VUgras{بالٹی}
\textsuperscript{(9)}۔

%%K t06-c1-04-1 p
4. --- \VUgras{باہرلے کمرے} \textsuperscript{(10)} وچ سانوں کجھ \VUgras{کھونٹیاں}
\textsuperscript{(13)} وکھائی دیندیاں نیں تے دو \VUgras{چھتریاں}
\textsuperscript{(11)}، جیہڑیاں \VUgras{چھتری دان} \textsuperscript{(12)} وچ نیں۔

%%K t06-c1-05-1 p
5. --- دروازے دے دوجے پاسے \VUgras{شیشے والی الماری} \textsuperscript{(14)} کھڑی
اے، جیدے وچ \VUgras{کپڑے} \textsuperscript{(15)} رکھے رہندے نیں۔

%%K t06-c1-06-1 p
6. --- \VUgras{نوکراݨی} \textsuperscript{(16)} \VUgras{پلنگ} \textsuperscript{(17)}
لا رہی اے، ایسے بہانے اوہدے وکھو وکھ حصے ویکھ لئیے۔ \VUgras{چارپائی}
\textsuperscript{(20)} اُتے اک چنگا \VUgras{کمانی والا گدا} \textsuperscript{(21)}
اے، جیدے اُتے جسٹین ہُݨے اُون دا \VUgras{گدا} \textsuperscript{(22)} وِچھائے گی۔ فیر
اوہ کرسیاں توں دو \VUgras{چادراں} \textsuperscript{(23)} تے \VUgras{کمبل}
\textsuperscript{(24)} چُکے گی۔ فیر اوہ پلنگ دے سرہاݨے \VUgras{لما سِرہاݨا}
\textsuperscript{(25)} تے \VUgras{سِرہاݨا} \textsuperscript{(26)} رکھے گی، جیہڑے ہالے
\VUgras{رضائی} \textsuperscript{(27)} نال \VUgras{پلنگ دے اگے دی دری}
\textsuperscript{(32)} اُتے پئے نیں۔ اوہ \VUgras{پردیاں} \textsuperscript{(19)} تے
\VUgras{چھتر} \textsuperscript{(18)} اُتے پراں دا جھاڑن پھیردی اے، تے اَدھا کم ہو
جاندا اے۔

%%K t06-c1-07-1 p
7. --- فیر اوہنوں \VUgras{پلنگ کول دی میز} \textsuperscript{(28)} تھوڑی سنوارنی
پوے گی، \VUgras{الارم گھڑی} \textsuperscript{(29)} وچ چابی بھرنی پوے گی،
\VUgras{رات دی بتی} \textsuperscript{(30)} تیار کرنی پوے گی، تے \VUgras{موم بتی}
\textsuperscript{(31)} ہٹا کے شمعدان صاف کرنا پوے گا۔

%%K t06-tit-2 sub
\VUsaut{5.0mm}
\VUpk{10.2pt}{40}{سیوݨ دی ٹوکری تے انگیٹھی دا سامان۔}
\VUsaut{3.0mm}

%%K t06-c1-08-1 p
8. --- \VUgras{سیوݨ دی ٹوکری} \textsuperscript{(34)}، جیہڑی \VUgras{تپائی}
\textsuperscript{(33)} کول رکھی اے، اوہنوں وی سنوارن دی بڑی لوڑ اے۔ اوہنوں
\VUgras{کڑھائی} \textsuperscript{(35)} تہہ کرنی پوے گی، \VUgras{سیوݨ دی میز}
\textsuperscript{(38)} وچ \VUgras{انگشتانہ}، \VUgras{دھاگہ} \textsuperscript{(36)}،
\VUgras{سُوئیاں} \textsuperscript{(37)} تے \VUgras{قینچی} \textsuperscript{(39)}
رکھݨیاں پوݨ گیاں، تے میز دا ڈھکݨ \VUgras{چابی} \textsuperscript{(40)} نال بند کرنا
پوے گا۔

%%K t06-c1-09-1 p
9. --- وچکار رکھی \VUgras{تپائی میز} \textsuperscript{(41)} اُتے جسٹین
\VUgras{صراحی} \textsuperscript{(42)} دا پاݨی بدلے گی۔ اوہ \VUgras{کھنڈ دان}
\textsuperscript{(43)} وچ \VUgras{کھنڈ} پائے گی، \VUgras{کھنڈ دی چمٹی}
\textsuperscript{(44)} صاف کرے گی تے \VUgras{کپڑیاں دا برش} \textsuperscript{(46)}
ہٹا دیوے گی۔

%%K t06-c1-10-1 p
10. --- کمرے دا سجا پاسا اوہنوں گھٹ کم دیوے گا، کیوں جے \VUgras{دن دا پلنگ}
\textsuperscript{(47)} تے اوہدی \VUgras{گدی} \textsuperscript{(48)} اپݨی تھاں اُتے
نیں، تے کیوں جے موسم ٹھنڈا نئیں، \VUgras{انگیٹھی} \textsuperscript{(49)} دے سامان
وچ کجھ ہلایا نئیں گیا۔ \VUgras{کہی} \textsuperscript{(50)}، \VUgras{چمٹا}
\textsuperscript{(51)}، \VUgras{لکڑ دے ٹیک} \textsuperscript{(55)} تے
\VUgras{سواہ دا پردہ} \textsuperscript{(56)} صاف نیں؛ \VUgras{اگ دا پردہ}
\textsuperscript{(54)} ہلایا نئیں گیا تے \VUgras{لکڑ دا صندوق} \textsuperscript{(52)}
\VUgras{بالݨ} \textsuperscript{(53)} نال بھریا اے۔

%%K t06-noto-1 noto
\VUnotes{143.2mm}{%
(*) Babo = نِکا بال؛ انگریزی \textit{baby}، فرانسیسی \textit{bébé}، اطالوی
\textit{bambino}۔%
}

%%K t06-noto-2 noto
\VUnotes{143.2mm}{%
(*) Lambrequino = فرانسیسی \textit{lambrequin}، ہسپانوی
\textit{lambrequines}، پرتگالی \textit{lambrequins}، تے جرمن تے انگریزی وچ وی
\textit{lambrequins} --- ارتھات جرمن، انگریزی، فرانسیسی، پرتگالی تے ہسپانوی،
سبھناں وچ اکو۔%
}

%%K t06-c1-11-1 p
11. --- اوہنوں ہالے \VUgras{انگیٹھی دی گھڑی} \textsuperscript{(57)}،
\VUgras{شمعدان} \textsuperscript{(58)} تے \VUgras{سجاوٹ دیاں طشتریاں}
\textsuperscript{(59)} جھاڑنیاں پوݨ گیاں، فیر \VUgras{شیشہ} \textsuperscript{(60)}
پونجھݨا پوے گا۔

%%K t06-c1-12-1 p
12. --- رہیا بال دا \VUgras{پنگھوڑا} \textsuperscript{(61)} (اوہ babo (*))، اوہ تے
پہلاں توں تیار اے۔

%%K t06-tit-3 sub
\VUsaut{5.0mm}
\VUpk{10.2pt}{40}{بیٹھک۔}
\VUsaut{3.0mm}

%%K t06-c2-01-1 p
1. --- ہُݨ ایہہ رہی بیٹھک۔ ایہہ بہت سوہݨا \VUgras{کمرہ} اے۔ سجے کونے دی انگیٹھی
اک نازک \VUgras{مورتی} \textsuperscript{(1)} تے دو سوہݨے \VUgras{پھُلداناں}
\textsuperscript{(3)} نال سجی اے، تے مخمل دی \VUgras{جھالر} \textsuperscript{(2)} (*)
نال ڈھکی اے۔

%%K t06-c2-02-1 p
2. --- انگیٹھی دیاں چنگیاں قالین نوں اگ نہ لا دیݨ، ایسے لئی اوہدے سامݨے تانبے دا
اک \VUgras{اگ دا جنگلا} \textsuperscript{(4)} رکھیا گیا اے۔

%%K t06-c2-03-1 p
3. --- اوہدے کول اک سجیلی \VUgras{لکھݨ دی میز} \textsuperscript{(5)} کھلی کھڑی اے۔
ایتھے ای \VUgras{گھر دی مالکݨ} \textsuperscript{(23)} اپݨے خط لکھدی اے۔

%%K t06-c2-04-1 p
4. --- کھڑکی اُتے \VUgras{جالی دار پردے} \textsuperscript{(6)} لگے نیں، جیہناں نوں
\VUgras{ڈوریاں} \textsuperscript{(8)} بنھی رکھدیاں نیں، تے اوہ ڈوریاں
\VUgras{کھونٹیاں} \textsuperscript{(7)} وچ اٹکیاں نیں، تے اوہناں اُتے بھارے کپڑے دے
پردے، جیہناں دے اُتے اک \VUgras{جھالر پٹی} \textsuperscript{(9)} اے۔

%%K t06-c2-05-1 p
5. --- کھڑکی دے طاق وچ اک \VUgras{گملا چوکی} \textsuperscript{(11)} وچ ہرا
\VUgras{بوٹا} \textsuperscript{(10)} اے، اک شاندار کھجور، جیدے پتّر تقریباً
\VUgras{تیل دے لیمپ} \textsuperscript{(12)} تیکر اپڑدے نیں۔

%%K t06-c2-06-1 p
6. --- \VUgras{لیمپ دی چھاں} \textsuperscript{(13)} ماریا نے بݨائی اے، اوہو ای موہ
لیݨ والی \VUgras{مُٹیار} \textsuperscript{(14)} جیہڑی \VUgras{پیانو}
\textsuperscript{(15)} اُتے بیٹھی اے۔ \VUgras{چوکی} \textsuperscript{(16)} اُتے سِدھی
بیٹھی اوہ اک دھن دی مشق کر رہی اے۔ اوہ بہت ہوشیار تے بہت سلیقے والی اے، جِویں
اوہدی \VUgras{سُر لپی دی الماری} \textsuperscript{(17)} توں پتہ لگدا اے: اوہدے وچ
پورا ترتیب اے۔

%%K t06-tit-4 sub
\VUsaut{5.0mm}
\VUpk{10.2pt}{40}{شرارتی۔}
\VUsaut{3.0mm}

%%K t06-c2-07-1 p
7. --- اوہدا بھرا پال اک \VUgras{کتاب} \textsuperscript{(19)} لبھ رہیا اے،
\VUgras{کتاباں دی الماری} \textsuperscript{(18)} وچ۔ اوہ ایہو جیہا \VUgras{مُنڈا}
\textsuperscript{(20)} اے جیہڑا شرارتی تے بالکل بے قابو اے۔ بہت اُچی رکھی کتاب تیکر
اپڑن لئی الماری توں لٹکدا ہویا اوہ اوہدے اُتے رکھی \VUgras{مورتی}
\textsuperscript{(21)} ڈیگ دیݨ دا خطرہ چُک رہیا اے۔

%%K t06-c2-08-1 p
8. --- ایہہ شرارت کرن لئی اوہ ایس گل دا فائدہ چُک رہیا اے پئی اوہدی ماں اک سہیلی
نال گلاں وچ لگی اے، اک \VUgras{بیبی} \textsuperscript{(22)} جیہڑی کجھ دن اوہناں نال
لنگھاؤݨ آئی اے۔ ماں \VUgras{صوفے} \textsuperscript{(24)} اُتے بیٹھی اے،
\VUgras{آرام کرسی} \textsuperscript{(25)} کول، تے اوہدی سہیلی اوس \VUgras{کرسی}
\textsuperscript{(27)} اُتے، جیدی سانوں نِری \VUgras{پِٹھ} \textsuperscript{(26)}
وکھائی دیندی اے۔

%%K t06-c2-09-1 p
9. --- کندھ اُتے ٹنگے وڈے \VUgras{فریم} \textsuperscript{(30)} وچ اک رشتے دار دی
\VUgras{تصویر} \textsuperscript{(29)} اے۔ میز اُتے پئی \VUgras{البم}
\textsuperscript{(32)} وچ باقی ٹبر دیاں \VUgras{تصویراں} \textsuperscript{(33)}
کٹھیاں نیں۔ نِکے ہُݨے تھوڑی دیر پہلاں اوہنوں پُٹھ پرت رہے سن، کیوں جے
\VUgras{کرسیاں} \textsuperscript{(31)} ہالے وی میز دے چوہاں پاسیں جمع نیں۔

%%K t06-c2-10-1 p
10. --- چینی مِٹی دا \VUgras{چینی پھُلدان} \textsuperscript{(34)}، جیہڑا
\VUgras{کاغذ کٹݨ دی چھُری} \textsuperscript{(35)} کول رکھیا اے، ماریا نوں اوہدے ناں
دے دن مِلیا سی، تے کمرے دا کونا بھردا \VUgras{تہہ ہوݨ والا پردہ}
\textsuperscript{(36)} اوہدے چاچا چین توں لیائے سن۔

%%K t06-c2-11-1 p
11. --- سوہݨیاں \VUgras{تصویراں} \textsuperscript{(37)} نال جگمگاندی، جیہڑیاں
\VUgras{کندھ} \textsuperscript{(42)} اُتے ٹنگیاں نیں، چھت دے \VUgras{جھاڑ فانوس}
\textsuperscript{(38)} نال روشن، جیدے \VUgras{بجلی دے لیمپ} \textsuperscript{(39)}
شام نوں اِنّے خوش چمکدے نیں، ایہہ بیٹھک بڑی سوہݨی لگدی اے۔ فرش اُتے وِچھیا کُونا
\VUgras{قالین} \textsuperscript{(40)} تے اِنّی کم دی \VUgras{بجلی دی گھنٹی}
\textsuperscript{(41)} اوہنوں پوری طرحاں آرام دیݨ والی بݨا دیندے نیں۔

%%K t06-tit-5 sub
\VUsaut{3.6mm}
\VUpk{10.2pt}{40}{کھاݨے دا کمرہ۔}
\VUsaut{3.0mm}

%%K t06-c3-01-1 p
1. --- کھاݨے دی گھنٹی وج چُکی اے۔ ماریا نوں چھڈ کے سارا ٹبر میز دے چوہاں پاسیں
جمع اے۔ کھاݨے دا کمرہ اک ودھیا مِٹی دے \VUgras{چولھے} \textsuperscript{(1)} نال
سوہݨا نیما اے، جیدی \VUgras{نالی} \textsuperscript{(2)} پتلی اے۔

%%K t06-c3-02-1 p
2. --- ایس کمرے وچ سامان بہت نئیں۔ کندھ نال، \VUgras{تپائی} \textsuperscript{(4)}
دے اُتے، اک سجاوٹی نِکا \VUgras{فوارہ} \textsuperscript{(3)} ٹنگیا اے۔ کول ای
\VUgras{سجاوٹی الماری} \textsuperscript{(5)} اے، جیدے اُتے ٹھنڈے کھاݨے، مِٹھیاں تے
اوہ برتن رکھے جاندے نیں جیہناں دی ایس ویلے لوڑ نئیں۔

%%K t06-c3-03-1 p
3. --- سو اُتلے خانے وچ سانوں اک \VUgras{بھُنیا کُکڑ} \textsuperscript{(7)} وکھائی
دیندا اے، اک \VUgras{مچھی} \textsuperscript{(8)} تے اک \VUgras{کٹورا}
\textsuperscript{(10)} \VUgras{پھلاں} \textsuperscript{(9)} دا۔ تھلے \VUgras{پنیر}
\textsuperscript{(11)}، \VUgras{روٹی} \textsuperscript{(12)} تے
\VUgras{تیل سرکے دا دان} \textsuperscript{(13)} اے، جیدے وچ سلاد سجاؤݨ دا سب کجھ
رہندا اے: تیل، سرکہ، \VUgras{لُوݨ} \textsuperscript{(14)} تے \VUgras{مرچ}
\textsuperscript{(15)}۔ چھیکر سب توں تھلے خانے وچ اک \VUgras{کیک} \textsuperscript{(17)}
اے، \VUgras{رائی دے دان} \textsuperscript{(16)} کول۔

%%K t06-c3-04-1 p
4. --- کھاݨے دے کمرے دی گھڑی، جیہنوں \VUgras{کندھ گھڑی} \textsuperscript{(20)}
آکھدے نیں، بہت انوکھی شکل دی اے۔ اوہ لکڑ دے اک چوکھٹے وچ گڈی اے، جیدے وچ اک
\VUgras{ہوا دباؤ ماپ} \textsuperscript{(19)} تے اک \VUgras{تاپ ماپ}
\textsuperscript{(18)} وی نیں۔

%%K t06-tit-6 sub
\VUsaut{4.6mm}
\VUpk{10.2pt}{40}{کھاݨا وَرتاؤݨا۔}
\VUsaut{3.0mm}

%%K t06-c3-05-1 p
5. --- کمرے دے چوہاں پاسیں کندھاں اُتے موم دی پالش کیتی بلوط دی
\VUgras{لکڑ دی چوکھٹ} \textsuperscript{(21)} لگی اے۔ کپڑے دا اک
\VUgras{دروازے دا پردہ} \textsuperscript{(22)} برامدے ول جاندے دروازے نوں چنگی طرحاں
ڈھک دیندا اے۔ کھاݨے مگروں ٹبر اوتھے ای کافی پیݨ جاوے گا۔ \VUgras{پیالے}
\textsuperscript{(24)} تے \VUgras{طشتریاں} \textsuperscript{(25)} پہلاں توں
\VUgras{چینی برتناں دی الماری} \textsuperscript{(23)} اُتے سجے نیں۔

%%K t06-c3-06-1 p
6. --- میز دے اُتے چھت توں اک \VUgras{لٹکدا لیمپ} \textsuperscript{(26)} بنھیا اے؛
اوہدی بہت تیز لو شام نوں سارے کھاݨے دے کمرے وچ بھر جاندی اے۔

%%K t06-c3-07-1 p
7. --- ہُݨ آپ \VUgras{کھاݨے دی میز} \textsuperscript{(27)} ویکھیے۔ ٹبر دے آوݨ اُتے
میز لگی ہوئی سی۔ \VUgras{میز پوش} \textsuperscript{(28)} اُتے ہر کھاݨ والے دی تھاں
رکھی سی اک \VUgras{پلیٹ} \textsuperscript{(33)}، اک \VUgras{رومال}
\textsuperscript{(29)}، اک \VUgras{چمچہ} \textsuperscript{(34)}، اک \VUgras{کانٹا}
\textsuperscript{(35)}، اک \VUgras{چھُری} \textsuperscript{(36)} تے اک
\VUgras{گلاس} \textsuperscript{(32)}۔ شراب دیاں \VUgras{بوتلاں} \textsuperscript{(30)}
تے پاݨی دی اک \VUgras{صراحی} \textsuperscript{(31)} وی سی۔

%%K t06-c3-08-1 p
8. --- \VUgras{نوکر} \textsuperscript{(39)} پہلاں \VUgras{سُوپ دا برتن}
\textsuperscript{(37)} لیایا، جیدے وچوں ہالے وی کجھ سُوپ توں بھاپ اُٹھ رہی اے۔ ہُݨ
اوہ پہلا \VUgras{کھاݨا} \textsuperscript{(38)} وَرتا رہیا اے، مانس دا۔ فیر اوہ اوہ
کھاݨے لیائے گا جیہناں دے ناں اسیں لے چُکے آں؛ تے چھیکر، \VUgras{مِٹھے} مگروں،
اوہ ایس گل دا فائدہ چُکے گا پئی مالک برامدے وچ چلے گئے نیں، تے میز خالی کر کے
کھاݨے دا کمرہ مُڑ ٹھیک کر دیوے گا۔

%%K t06-c3-08-2 p
\textit{ٹُوک۔} --- \textit{کھاݨے نال جُڑے روز دے لفظ ایتھے نِرے موٹے طور اُتے
دِتے گئے نیں۔ ایہہ لسٹ « ہوٹل » (نمبر 13) تے « منڈی » (نمبر 15) دے دو نقشیاں وچ
پوری کیتی جاوے گی۔}

%%K t06-tit-7 sub
\VUsaut{4.6mm}
\VUpk{10.2pt}{40}{رسوئی۔}
\VUsaut{3.0mm}

%%K t06-c4-01-1 p
1. --- اَدھ کھلے دروازے وچوں جیس رسوئی اُتے اسیں نظر مارنے آں، اوتھے اک سوہݨا
گڑبڑ جھالا لبھدا اے۔ \VUgras{چولھے} \textsuperscript{(1)} دے سامݨے، جِتھے
\VUgras{اگ} \textsuperscript{(3)} چٹک رہی اے، \VUgras{بھُنݨ دا جنتر}
\textsuperscript{(5)} لگا اے۔ اک ودھیا \VUgras{بھُنیا مانس} \textsuperscript{(7)}
\VUgras{سیخ} \textsuperscript{(6)} اُتے اے تے پکݨ نوں اے۔

%%K t06-c4-02-1 p
2. --- \VUgras{دھوکݨی} \textsuperscript{(8)} تے \VUgras{کوئلے دی کہی}
\textsuperscript{(9)}، جیہڑیاں ہُݨے کم آئیاں نیں، \VUgras{لکڑاں}
\textsuperscript{(10)} سمیت زمین اُتے ای چھڈ دِتیاں گئیاں نیں۔

%%K t06-c4-03-1 p
3. --- انگیٹھی دے اُتے دا حصہ سنوریا ہویا اے: \VUgras{لالٹین}
\textsuperscript{(11)}، \VUgras{دیاسلائی دان} \textsuperscript{(13)} جیدے وچ
\VUgras{دیاسلائیاں} \textsuperscript{(12)} نیں، \VUgras{شمعدان} \textsuperscript{(14)}
تے \VUgras{ہتھ دا شمعدان} \textsuperscript{(15)} اپݨی \VUgras{موم بتی}
\textsuperscript{(16)} سمیت — سب اپݨی تھاں نیں۔ پر وڈی \VUgras{کوئلے دی ٹوکری}
\textsuperscript{(18)}، جیہڑی اوس \VUgras{کوئلے} \textsuperscript{(17)} نال بھری اے
جیہڑا \VUgras{رسوئݨ} \textsuperscript{(23)} ہُݨے تہہ خانے توں لیائی اے، رسوئی دے
وچکار ای پئی رہ گئی اے۔

%%K t06-c4-04-1 p
4. --- سچ تے ایہہ اے پئی وچاری نوں چھیتی کرنی اے، جے دوپہر دا کھاݨا ویلے سِر
تیار کرنا اے۔ ایس ویلے اوہ \VUgras{پکاؤݨ دے چولھے} \textsuperscript{(19)} اُتے اک
\VUgras{رَسا} \textsuperscript{(24)} چلا رہی اے، جیہڑا سڑن نہ پاوے۔

%%K t06-c4-05-1 p
5. --- \VUgras{تندور} \textsuperscript{(20)} وچ اک کیک پک رہیا اے، جیہڑا اوہنے
نِکیاں دی شام دی چا لئی بݨایا اے۔ چولھے دے اُتے \VUgras{کڑچھی}
\textsuperscript{(21)} تے \VUgras{جھرنا} \textsuperscript{(22)} ٹنگے نیں۔

%%K t06-tit-8 sub
\VUsaut{4.0mm}
\VUpk{10.2pt}{40}{رسوئی دے برتن۔}
\VUsaut{3.0mm}

%%K t06-c4-06-1 p
6. --- کمرے دے پچھے ٹنگیا \VUgras{برتناں دا سیٹ} \textsuperscript{(25)} بہت صاف اے۔
تانبے دیاں سوہݨیاں \VUgras{پتیلیاں} \textsuperscript{(26)}، \VUgras{دُدھ دا جگ}
\textsuperscript{(27)}، \VUgras{چھاننی} \textsuperscript{(28)} تے \VUgras{کدوکش}
\textsuperscript{(29)} دھیان نال مانجے گئے نیں۔

%%K t06-c4-07-1 p
7. --- \VUgras{لُوݨ دان} \textsuperscript{(30)} نِکی الماری اُتے \VUgras{چا دان}
\textsuperscript{(31)} تے \VUgras{تیل دی وڈی بوتل} \textsuperscript{(32)} کول رکھیا
اے۔ \VUgras{دراز} \textsuperscript{(33)} وچ شاید چمچے کانٹے تے رسوئی دیاں چھُریاں
رہندیاں نیں۔

%%K t06-c4-08-1 p
8. --- \VUgras{جھاڑو} \textsuperscript{(34)} تے \VUgras{پراں دا جھاڑن}
\textsuperscript{(35)} اک کونے وچ نیں۔ رسوئݨ ہُݨے \VUgras{کٹݨ دا کُندا}
\textsuperscript{(36)} کم وچ لیا چُکی اے؛ اوہنے اوہنوں کھڑکی کول چھڈ دِتا اے۔

%%K t06-c4-09-1 p
9. --- اک \VUgras{ہتھ دا تولیہ} \textsuperscript{(37)} کندھ اُتے ٹنگیا اے،
\VUgras{قیف} \textsuperscript{(38)} کول۔ رسوئی دے ایس حصے وچ \VUgras{سیخ جالی}
\textsuperscript{(39)}، \VUgras{کٹار} \textsuperscript{(40)} تے
\VUgras{کٹݨ دا پٹرا} \textsuperscript{(41)} رکھے جاندے نیں۔ فیر کٹار دے اُتے، اک
\VUgras{طاق} \textsuperscript{(42)} اُتے، \VUgras{ہانڈیاں} \textsuperscript{(43)}،
\VUgras{سالن دی دیگچی} \textsuperscript{(44)} اپݨے \VUgras{ڈھکݨ} \textsuperscript{(45)}
سمیت، تے \VUgras{کڑاہا} \textsuperscript{(46)} نیں۔ \VUgras{تلݨ دی کڑاہی}
\textsuperscript{(47)} کول ای ٹنگی اے۔

%%K t06-c4-10-1 p
10. --- ایس کونے وچ نِری تانبے دی \VUgras{ٹُوٹی} \textsuperscript{(48)} چمک پاندی
اے۔ \VUgras{نالی چوکی} \textsuperscript{(49)}، جیدے اُتے وڈا \VUgras{گھڑا}
\textsuperscript{(50)} رکھیا اے، اپݨی نِکی الماری نال بہت کم دی ہووے گی، جِتھے
\VUgras{سرکے دا پیپا} \textsuperscript{(51)} اپݨی \VUgras{ٹُوٹی}
\textsuperscript{(52)} سمیت، \VUgras{کوڑے دا صندوق} \textsuperscript{(53)} تے
\VUgras{گندے برتن} \textsuperscript{(54)} رکھے جاندے نیں۔

%%K t06-tit-9 sub
\VUsaut{4.0mm}
\VUpk{10.2pt}{40}{رسوئی وچ رکھیاں ہور چیزاں۔}
\VUsaut{3.0mm}

%%K t06-c4-11-1 p
11. --- اوہ وڈا \VUgras{ٹب} \textsuperscript{(55)} جیدے وچ \VUgras{ترکاریاں}
\textsuperscript{(58)} دھوتیاں جاندیاں نیں، تے ڈھالے لوہے دا \VUgras{کڑاہا}
\textsuperscript{(56)}، جیہڑا \VUgras{اِٹ دے فرش} \textsuperscript{(57)} اُتے رکھیا
اے، شاید اوسے الماری دے نیں۔

%%K t06-c4-12-1 p
12. --- رسوئݨ کول چولھیاں دی گھاٹ نئیں، کیوں جے ایہہ رہیا میز دے کونے اُتے اک
\VUgras{گیس دا چولھا} \textsuperscript{(59)}، جیدے اُتے اک نِکی پتیلی وچ پاݨی گرم ہو
رہیا اے۔ اوہدے وچوں نکلدی \VUgras{بھاپ} \textsuperscript{(60)} دسدی اے پئی اوہ اُبل
رہیا ہووے گا۔ ایہہ پاݨی شاید کافی لئی اے، کیوں جے میز دے دوجے سرے اُتے ایہہ رہے
\VUgras{کافی دا برتن} \textsuperscript{(64)} تے \VUgras{کافی دی چکی}
\textsuperscript{(63)}۔

%%K t06-c4-13-1 p
13. --- گیس دا چولھا \VUgras{گیس دی ٹُوٹی} \textsuperscript{(62)} نال جُڑیا اے، اک
لمی \VUgras{ربڑ دی نالی} \textsuperscript{(61)} دے ذریعے۔

%%K t06-c4-14-1 p
14. --- \VUgras{دیہاڑی دی مددگار} \textsuperscript{(66)}، جیہڑی رسوئݨ دی مدد نوں
آندی اے، رات دے کھاݨے دیاں ترکاریاں تیار کر رہی اے۔ چھِل کے تے دھو کے اوہ اوہناں
نوں \VUgras{برتن} \textsuperscript{(65)} وچ رکھ دیوے گی، جدوں تیکر پکاؤݨ دا ویلا نہ
آوے۔

%%K t06-tit-10 sub
\VUsaut{4.0mm}
{\centering\fontsize{10.2pt}{10.2pt}\selectfont\textls[40]{\scshape غسل خانہ۔} \textit{(نہاؤݨ دا کمرہ۔)}\par}
\VUsaut{3.0mm}

%%K t06-c5-01-1 p
1. --- گھر دے مکھ کمرے جاݨ لیݨ لئی ہُݨ اکو جھاتی باقی اے: غسل خانے دا سامان
ویکھݨا۔ ایس وچ دیر نہ لگے گی، کیوں جے اوہ وڈا نئیں، تے اسیں جھٹ ویکھ لواں گے پئی
\VUgras{نہاؤݨ دے ٹب} \textsuperscript{(1)} نال اک چنگا \VUgras{پاݨی گرم کرن دا
جنتر} \textsuperscript{(2)} اے، پئی \VUgras{صابن دانی} \textsuperscript{(3)} نہاؤݨ
والے دے ہتھ دی پہنچ وچ اے، تے پئی \VUgras{تولیے دے ڈنڈے} \textsuperscript{(5)} نال
\VUgras{تولیے} \textsuperscript{(4)} سُکاؤݨا سوکھا ہوندا ہووے گا۔

%%K t06-c5-02-1 p
2. --- ایس کمرے دیاں کندھاں \VUgras{چمکدیاں ٹائلاں} \textsuperscript{(6)} نال ڈھکیاں
نیں، جیدے نال اک \VUgras{فوارہ} \textsuperscript{(7)} لاؤݨا ممکن ہویا، بغیر ایس ڈر
دے پئی اوہنوں ورتدے ویلے چھِٹے اوہناں نوں وگاڑ دیݨ۔ پیر دھوݨ دا جست دا اک نِکا
\VUgras{تسلہ} \textsuperscript{(8)} سامان پورا کر دیندا اے۔
\VUsaut{6.0mm}
\VUfilet{22mm}
